\section{Further research questions and exact-identity programme}
\label{sec:questions}

The following questions begin where the proved formulas stop. They
are proposed research directions, not assertions disguised as
conjectures. Their intended outputs are exact identities, finite
coordinate descriptions, or proofs of specific candidate equalities.

\subsection{Independent slopes beyond depth three}

\paragraph{Q1. The first genuinely new ordered coefficients.}
At depth four the finite part of a general ray uses the linear jet
of $R_3$ and the quadratic jet of $R_2$, in addition to ordinary
Stieltjes data. Determine a complete finite basis of these
coefficients modulo the exact stuffle relations, shift identities,
and the known diagonal Gamma constraints. Then write a formula
analogous to \eqref{ord:threeFP} for all transverse quadruples of
slopes. The first objective is a proved spanning description; any
claim of arithmetic independence would require a separate argument.

\paragraph{Q2. Classify the directions with ordinary-zeta closure.}
The family with equal trailing slopes has ordinary Hurwitz closure
at every depth. At depth three the coefficient of the ordering
coordinate is $(c_2-c_3)/c_1$. At depth four and beyond, determine
the algebraic loci on which all genuinely nonsymmetric coefficient
contributions cancel. One can first work in the formal quotient by
proved identities, where the question becomes an exact finite
linear-algebra problem. Accidental numerical relations at a specific
shift must be distinguished from identities valid as functions of $a$.

\paragraph{Q3. Several blocks of equal slopes.}
Replace $(p,q,\ldots,q)$ by a sequence of several constant blocks.
Does each fixed number of blocks admit a generating function involving
a finite collection of lower-depth regular germs, uniformly in the
block lengths? The all-equal and one-distinguished-slope cases suggest
an induction by ordered blocks, but the crossing terms need explicit
control. A useful deliverable would be an exact two-block formula at
all depths and a proof of its coordinate requirements.

\subsection{Harmonic, polylogarithmic, and Gamma relations}

\paragraph{Q4. Rational-shift distribution of the ordering coordinate.}
The shift relation \eqref{harm:shift} is exact. The existing
root-of-unity filter in the incoming endpoint report already gives
finite coloured descriptions of rational-shift ordered sums and
their fixed spectral derivatives \cite{repoendpoint}. The harder
question is which character-weighted combinations of $\eta(r/q)$
eliminate the residual coloured ordering coefficients and reduce
entirely to ordinary Hurwitz or Dirichlet $L$-derivatives. Determine
the exact elimination spaces first at conductors $2,3,4,6$, then
seek a conductor-uniform rule. The inclusive harmonic literature
\cite{Kargin} is a starting point; the strict/inclusive conversion
must remain explicit.

\paragraph{Q5. Normalized parameter primitives.}
Equation~\eqref{harm:derivative} reduces $\eta'(a)$ to absolutely
convergent ordered sums. Seek a useful normalized antiderivative in
the shift parameter, expressed through multiple Gamma or Barnes-type
functions and spectral zeta derivatives. A candidate must specify
its additive constant and prove differentiation back to $\eta$.
The whole higher harmonic hierarchy is already represented by
\eqref{harm:allcoefficients}--\eqref{harm:allconversion}; seek a
parameter-primitive relation compatible with that hierarchy and with
the exact shift recurrence. Any regularized integral must carry its
finite normalization explicitly.

\paragraph{Q6. Logarithmic decorations of convergent primitive families.}
One Euler integration raises the outer order from one to two, so
$Z_{r+1}(2+u_0,1+u_1,\ldots,1+u_r;a)$ is holomorphic near the origin.
For this ordinary primitive the endpoint limit and every fixed
spectral derivative commute by domination; there is no unexplained
regularization defect. The incoming endpoint work already evaluates
the unweighted integer-shift primitive hierarchy
\cite{repoendpoint}. The next target is to classify which independent
logarithmic decorations, obtained by differentiating the separate
indices, still reduce to ordinary Gamma, zeta and Stieltjes jets.
Begin with the convergent depth-two derivatives appearing in
\eqref{harm:derivative}, derive their exact integral relations, and
compare them with the height-one Gamma generating relations.

\subsection{Tornheim derivatives and higher coincident products}

\paragraph{Q7. Evaluate the remaining cubic Tornheim derivative.}
Can $\tau'''(0)$ be reduced to a specified algebra generated by
$\gamma_m$, ordinary $\zeta^{(j)}(n)$, logarithms, and perhaps Barnes
Gamma constants? Formula (C9) turns this into an equivalent exact
evaluation question for the coordinate-normalized digamma cube.
The independent Mellin representation permits systematic analytic
manipulations and high-precision exploration, but an integer relation
would remain a candidate until a functional identity proves it.
Existing Tornheim derivative work and the recent Kronecker-limit
literature should be checked before any priority claim
\cite{BorweinDilcher,BaileyBorwein,Onodera2021,SS2026}.

\paragraph{Q8. Minimal coordinates for all cubic moments.}
Formula (C4) is a complete generating prescription. It does not yet
classify how many independent completed Tornheim jets are needed at
each total Stieltjes index $m_1+m_2+m_3$. Determine exact reductions
under permutation symmetry, integration by parts with the stated
finite-part convention, and Tornheim functional equations. The
first useful target is a table for total indices through three,
with proofs and a transparent distinction between a spanning family
and an arithmetically minimal one.

\paragraph{Q9. Four and more coincident factors.}
Extend the cubic Fourier kernel to the zero-frequency sums arising
from four Hurwitz factors. The endpoint resonances can be organized
by subsets of singular factors and their smooth Taylor degrees;
the genuinely new work is the global multiple-sum coordinate and its
functional reductions. Begin with a fixed-coordinate formula for
$\FP\int_0^1\psi(x)^4\,dx$, keeping all spectral-counterterm finite
contributions. The natural higher-depth function should be defined
independently, as (C6) defines $\tau$ here.

\subsection{Singular geometry and the Gaussian programme}

\paragraph{Q10. Restrictions lying inside a polar hyperplane.}
Theorem~\ref{curve:main} handles every analytic curve whose prefixes
are not identically zero. A curve contained in $P_j=0$ needs an
additional prescription. Compare residue-first, regular-germ-first,
and finite-part-in-a-normal-variable constructions, and derive their
exact differences. Decide which prescriptions respect the stuffle
product and how those choices propagate under changes of coordinates.
This should be formulated as a comparison theorem; no canonical
value follows just from substituting into a rational slope formula.

\paragraph{Q11. A functional proof of the short Gaussian candidates.}
The current $S_6$ and revised $S_8$ candidates require identities
between specified polylogarithmic periods. A productive next step
would be an exact functional relation, differential equation with
enough boundary data, or a finite symbolic certificate in a rigorously
defined iterated-integral algebra. Existing restricted separator
results constrain particular ansatz classes; they do not establish
general period independence. The ordered Hurwitz identities in this
article may supply useful coefficient relations, but no demonstrated
bridge presently converts them into a proof of either Gaussian
candidate. That bridge itself should be stated as a concrete research
target rather than assumed.
